\documentclass{article}
%%%%%%%%%%%%%
% Loads packages
%%%%%%%%%%%%%
\usepackage[table]{xcolor}
\usepackage[utf8]{inputenc}
\usepackage[colorlinks=true,linkcolor=blue]{hyperref}
\usepackage{geometry} %package needed to set margins
\usepackage{fancyhdr}
\usepackage{graphicx}
\usepackage{amsmath}
\usepackage{amsthm}
\usepackage{mdframed}
\usepackage{tikz}
\usepackage{amsfonts}

\usepackage{listings}% http://ctan.org/pkg/listings
\lstset{
  basicstyle=\ttfamily,
  mathescape
}


\pagestyle{fancy}
\fancyhf{}
\chead{\textbf{Homework 5}}
\lhead{Math 213, Spring 2026}
\rhead{Due Wednesday, 3/4 at 9:00am}

%%%%%%%%%%%%%
% Sets margins
%%%%%%%%%%%%%
\newgeometry{left=1.5in,right=1in,top=1in,bottom=1in}
\setlength\headsep{3pt}

%%%%%%%%%%%%%
% Creates problem and solution environments
%%%%%%%%%%%%%

% Solution Environment
\newenvironment{solution}{\begin{proof}[Solution]}{\end{proof}}

% Problem Environment
\newenvironment{problem}[1]
    {\begin{mdframed}[default]
    \textbf{Problem #1:}
    }
    {\end{mdframed}
    }
    
%%%%%%%%%%%
% Custom Commands
%%%%%%%%%%%
\newcommand{\gOne}{\cellcolor{green!50!white} 1}
\newcommand{\rZero}{\cellcolor{red!50!white} 0}

\begin{document}

\begin{problem}{\S 6.2: 2}
Show that if there are 30 students in a class, then at least two have last names that begin with the same letter.
\end{problem}

\begin{problem}{\S 6.2: 6}
Let $d$ be a positive integer. Show that among any group of $d+1$ (not necessarily consecutive) integers there are two with exactly the same remainder when they are divided by $d$.
\end{problem}

\begin{problem}{\S 6.2: 8}
Show that if $f$ is a function from $S$ to $T$, where $S$ and $T$ are finite sets with $|S| > |T|$, then there are elements $s_1$ and $s_2$ in $S$ such that $f(s_1) = f(s_2)$, or in other words, $f$ is not one-to-one.
\end{problem}

\begin{problem}{\S 6.2: 14}
\begin{enumerate}
    \item[(a)] Show that if seven integers are selected from the first 10 positive integers, there must be at least two pairs of these integers with the sum 11.
    \item[(b)] Is the conclusion in (a) true if six integers are selected rather than seven?
\end{enumerate}
\end{problem}

\begin{problem}{\S 6.3 - 16}
How many subsets with an odd number of elements does a set with 10 elements have?
\end{problem}

\begin{problem}{\S 6.3 - 20}
How many bit strings of length 10 have
\begin{enumerate}
    \item[(a)] exactly three ``0''s?
    \item[(b)] more ``0''s than ``1''s?
    \item[(c)] at least seven ``1''s?
    \item[(d)] at least three ``1''s?
\end{enumerate}
\end{problem}

\begin{problem}{\S 6.3 - 24}
How many ways are there for ten women and six men to stand in a line so that no two men stand next to each other?
\end{problem}

\begin{problem}{\S 6.3 - 42}
Find a formula for the number of ways to seat $r$ of $n$ people around a circular table, where seatings are considered the same if every person has the same two neighbors without regard to which side these neighbors are sitting on.
\end{problem}

\begin{problem}{\S 6.4 - 10}
Give a formula for the coefficient of $x^k$ in the expansion of $(x+1/x)^{100}$, where $k$ is an integer.
\end{problem}

\begin{problem}{\S 6.4 - 12}
The row of Pascal's triangle containing the binomia coefficients ${ 10 \choose k}$, for $0 \leq k \leq 10$, is:
\[1~~10~~45~~120~~210~~252~~210~~120~~45~~10~~1 \]
Use Pascal's identity to produce the row immediately following this row in Pascal's triangle.
\end{problem}

\begin{problem}{\S 6.4 - 22}
Prove the identity ${n \choose r}{r \choose k} = {n \choose k}{n-k \choose r-k}$, whenever $n, r,$ and $k$ are non-negative integers with $r \leq n$ and $k \leq r$,
\begin{enumerate}
    \item[(a)] using a combinatorial argument.
    \item[(b)] using an argument based on the formula for the number of $r$-combinations of a set with $n$ elements.
\end{enumerate}
\end{problem}

\begin{problem}{\S 6.4 - 27(a)}
Prove the \textbf{hockeystick identity}
\[ \sum_{k=0}^r {n+k \choose k} = {n + r + 1 \choose r } \]
whenever $n$ and $r$ are positive integers, using a combinatorial argument.
\end{problem}

\begin{problem}{\S 6.5 - 10(a,c,d)}
A croissant shop has plain croissants, cherry croissants, chocolate croissants, almond croissants, apple croissants, and broccoli croissants. How many ways are there to choose
\begin{enumerate}
    \item[(a)] a dozen croissants?
    \item[(c)] two dozen croissants with at least two of each kind?
    \item[(d)] two dozen croissants with no more than two broccoli croissants?
\end{enumerate}
\end{problem}


\begin{problem}{\S 6.5 - 20}
How many solutions are there to the inequality
\[ x_1 + x_2 + x_3 \leq 11, \]
where $x_1, x_2,$ and $x_3$ are non-negative integers?
\end{problem}

\begin{problem}{\S 6.5 - 26}
How many positive integers less than $1,000,000$ have exactly one digit equal to $9$ and have a sum of digits equal to 13?
\end{problem}

\begin{problem}{\S 6.5 - 46}
A shelf holds 12 books in a row. How many ways are there to choose five books so that no two adjacent books are chosen?
\end{problem}

\end{document}